\begin{definition} \label{an1:def:metric-space}
  A \textbf{metric} on $X$ is a function $d:X\times X\to[0,\infty)$ such that, for all $p,q,z\in X$,
  \[
    d(p,q)=0\iff p=q,\qquad d(p,q)=d(q,p),\qquad d(p,z)\leq d(p,q)+d(q,z).
  \]
  The last condition is the \textbf{triangle inequality}. The pair $(X,d)$ is a \textbf{metric space}.
\end{definition}

% Lean source: Textbooks.MathematicalAnalysis1.Chapter01.EuclideanSpaces.euclidean_metric
\begin{proposition} \label{an1:thm:standard-metric-on-Rn}
  For $n\in\N$, the function $d(x,y)=\|x-y\|$ is a metric on $\R^{n}$.

  \begin{lean}
    theorem euclidean_metric (n : ℕ) :
        (∀ x y : EuclideanSpace ℝ (Fin n), 0 ≤ ‖x - y‖) ∧
        (∀ x y : EuclideanSpace ℝ (Fin n), ‖x - y‖ = 0 ↔ x = y) ∧
        (∀ x y : EuclideanSpace ℝ (Fin n), ‖x - y‖ = ‖y - x‖) ∧
        (∀ x y z : EuclideanSpace ℝ (Fin n),
          ‖x - z‖ ≤ ‖x - y‖ + ‖y - z‖)
  \end{lean}
\end{proposition}

\begin{proof}
  The norm is nonnegative. It vanishes exactly at the zero vector, so $\|x-y\|=0$ precisely when $x=y$.

  \begin{lean}
    refine ⟨fun x y => norm_nonneg _, ?_, ?_, ?_⟩

    · intro x y
      rw [norm_eq_zero, sub_eq_zero]
  \end{lean}

  Reversing a difference changes only its sign and leaves its norm unchanged.

  \begin{lean}
    · intro x y
      exact norm_sub_rev x y
  \end{lean}

  Write $x-z=(x-y)+(y-z)$ and apply the Euclidean norm triangle inequality.

  \begin{lean}
    · intro x y z
      calc
        ‖x - z‖ = ‖(x - y) + (y - z)‖ := by
          rw [sub_add_sub_cancel]
        _ ≤ ‖x - y‖ + ‖y - z‖ := norm_add_le _ _
  \end{lean}

\end{proof}

Euclidean spaces use this metric, including in dimension zero.

\begin{definition} \label{an1:def:neighborhood-open-ball}
  For $p\in X$ and $r>0$, the \textbf{open ball} is
  \[
    B(p,r)=\{q\in X:d(q,p)<r\}.
  \]
  The \textbf{punctured ball} is $B(p,r)\setminus\{p\}=\{q\in X:0<d(q,p)<r\}$.
  A \textbf{neighborhood} of $p$ is any set containing an open ball $B(p,r)$ with $r>0$. A neighborhood need not itself be open or a ball.
\end{definition}

\begin{definition} \label{an1:def:limit-interior-isolated-points}
  A point $p$ is a \textbf{limit point} of $E$ if every positive-radius punctured ball about $p$ meets $E$.
  The set of limit points, or \textbf{derived set}, is denoted by $E'$.
  A point is \textbf{interior} if some positive-radius ball about it lies in $E$; the set of interior points is denoted by $E^\circ$.
  An \textbf{isolated point} of $E$ belongs to $E$ and is its only point in some ball. We write $\operatorname{Iso}(E)$ for the set of isolated points.
  \begin{lean}
    def isolatedPoints (E : Set X) : Set X :=
      {p | ∃ r : ℝ, 0 < r ∧ E ∩ ball p r = {p}}
  \end{lean}
\end{definition}

% Lean source: Textbooks.MathematicalAnalysis1.Chapter01.MetricSpaces.mem_derivedSet_iff
\begin{lemma} \label{an1:lem:mem-derivedSet-iff}
  $p\in E^{\prime}$ if and only if every $r>0$ admits $q\in E$ with $q\ne p$ and $d(q,p)<r$.

  \begin{lean}
    theorem mem_derivedSet_iff (E : Set X) (p : X) :
        p ∈ derivedSet E ↔ ∀ r : ℝ, 0 < r → ∃ q ∈ E, q ≠ p ∧ dist q p < r
  \end{lean}
\end{lemma}

\begin{proof}
  Every positive-radius ball is a neighborhood, so the neighborhood condition supplies a distinct point in that ball.

  \begin{lean}
    rw [mem_derivedSet, accPt_iff_nhds]
    constructor

    · intro h r hr
      obtain ⟨q, ⟨hq, hE⟩, hne⟩ := h (ball p r) (Metric.ball_mem_nhds p hr)
      exact ⟨q, hE, hne, hq⟩
  \end{lean}

  Conversely, a neighborhood contains a positive-radius ball. A distinct point of the set in that ball also belongs to the neighborhood.

  \begin{lean}
    · intro h U hU
      obtain ⟨r, hr, hsub⟩ := Metric.mem_nhds_iff.mp hU
      obtain ⟨q, hE, hne, hq⟩ := h r hr
      exact ⟨q, ⟨hsub hq, hE⟩, hne⟩
  \end{lean}
\end{proof}

\begin{definition} \label{an1:def:open-closed-sets}
  A set $E$ is \textbf{open} if every point of $E$ is interior, and \textbf{closed} if its complement is open.
\end{definition}

\begin{definition} \label{an1:def:closure}
  The \textbf{closure} $\overline E$ consists of points whose every positive-radius ball meets $E$.
\end{definition}

\begin{definition} \label{an1:def:bounded-perfect-dense-sets}
  A set $E$ is \textbf{bounded} if there is a common upper bound for distances between its points.
  Equivalently, when $X$ is nonempty, $E$ lies in a positive-radius ball, as proved below. Subsets of the empty space are bounded.
  A set is \textbf{perfect} if it is closed and each of its points is a limit point.
  A set is \textbf{dense} if $\overline E=X$.
\end{definition}

% Lean source: Textbooks.MathematicalAnalysis1.Chapter01.MetricSpaces.isBounded_iff
\begin{lemma} \label{an1:lem:isBounded-iff}
  $E$ is bounded if and only if, whenever $X$ is nonempty, some $p\in X$ and $r>0$ satisfy $E\subseteq B(p,r)$.

  \begin{lean}
    theorem isBounded_iff (E : Set X) :
        Bornology.IsBounded E ↔ (Nonempty X → ∃ p : X, ∃ r > 0, E ⊆ ball p r)
  \end{lean}
\end{lemma}

\begin{proof}
  Suppose all pairwise distances in $E$ are at most $C$. If $E$ is empty, any center in $X$ and radius $1$ suffice. Otherwise choose $p\in E$ and radius $\max(C,0)+1$.

  \begin{lean}
    rw [Metric.isBounded_iff]
    constructor

    · rintro ⟨C, hC⟩ hX
      rcases E.eq_empty_or_nonempty with rfl | ⟨p, hp⟩

      · obtain ⟨p⟩ := hX
        exact ⟨p, 1, zero_lt_one, empty_subset _⟩

      · refine ⟨p, max C 0 + 1, by linarith [le_max_right C 0], ?_⟩
        intro q hq
        exact (hC hq hp).trans_lt (by linarith [le_max_left C 0])
  \end{lean}

  Conversely, two points of a ball of radius $r$ are at distance at most $2r$, by the triangle inequality. If $X$ is empty there are no points to check, and zero is a bound.

  \begin{lean}
    · intro h
      by_cases hX : Nonempty X

      · obtain ⟨p, r, _, hs⟩ := h hX
        refine ⟨r + r, ?_⟩
        intro x hx y hy
        have hx' : dist x p < r := hs hx
        have hy' : dist p y < r := by simpa only [mem_ball, dist_comm] using hs hy
        exact (dist_triangle x p y).trans (add_le_add hx'.le hy'.le)

      · refine ⟨0, ?_⟩
        intro x
        exact (hX ⟨x⟩).elim
  \end{lean}
\end{proof}

% Lean source: Textbooks.MathematicalAnalysis1.Chapter01.MetricSpaces.ball_open
\begin{proposition} \label{an1:thm:neighborhood-is-open}
  Every ball $B(p,r)$ is open. This also holds for nonpositive $r$, when the ball is empty.

  \begin{lean}
    theorem ball_open (p : X) (r : ℝ) : IsOpen (ball p r)
  \end{lean}
\end{proposition}

\begin{proof}
  Suppose $q\in B(p,r)$. The remaining distance $\delta=r-d(q,p)$ is positive.

  \begin{lean}
    apply Metric.isOpen_iff.mpr
    intro q hq
    refine ⟨r - dist q p, sub_pos.mpr hq, ?_⟩
  \end{lean}

  If $d(s,q)<\delta$, the triangle inequality gives $d(s,p)\leq d(s,q)+d(q,p)<r$. Thus the ball of radius $\delta$ about $q$ stays inside $B(p,r)$.

  \begin{lean}
    intro s hs
    change dist s p < r

    have ht := dist_triangle s q p

    change dist s q < r - dist q p at hs
    linarith
  \end{lean}

\end{proof}

% Lean source: Textbooks.MathematicalAnalysis1.Chapter01.MetricSpaces.finite_positive_lower_bound
\begin{lemma} \label{an1:lem:finite-positive-lower-bound}
  For a finite index set $S$ and positive numbers $r_i$ for $i\in S$, there is $\delta>0$ with $\delta\leq r_i$ for every $i\in S$.

  \begin{lean}
    theorem finite_positive_lower_bound {ι : Type*} (s : Finset ι)
        (r : ι → ℝ) (hr : ∀ i ∈ s, 0 < r i) :
        ∃ δ > 0, ∀ i ∈ s, δ ≤ r i
  \end{lean}
\end{lemma}

\begin{proof}
  Use induction on $S$. For the empty set, $\delta=1$ works. On adjoining $i$, take the smaller of $r_i$ and the positive bound already obtained for the remaining indices. It stays positive and is below every required radius.

  \begin{lean}
    classical

    induction s using Finset.induction_on with
    | empty => exact ⟨1, zero_lt_one, by simp only [Finset.notMem_empty, false_implies, implies_true]⟩
    | @insert i s hi ih =>
      obtain ⟨δ, hδ, hle⟩ := ih (fun j hj => hr j (Finset.mem_insert_of_mem hj))

      refine ⟨min (r i) δ, lt_min (hr i (Finset.mem_insert_self i s)) hδ, ?_⟩
      intro j hj
      rcases Finset.mem_insert.mp hj with rfl | hj

      · exact min_le_left _ _
      · exact (min_le_right _ _).trans (hle j hj)
  \end{lean}

\end{proof}

% Lean source: Textbooks.MathematicalAnalysis1.Chapter01.MetricSpaces.limit_point_infinite
\begin{proposition} \label{an1:thm:limit-point-neighborhood-contains-infinite-elements}
  If $p\in E^{\prime}$ and $r>0$, then $E\cap B(p,r)$ is infinite.

  \begin{lean}
    theorem limit_point_infinite (E : Set X) (p : X)
        (hp : p ∈ derivedSet E) (r : ℝ) (hr : 0 < r) :
        (E ∩ ball p r).Infinite
  \end{lean}
\end{proposition}

\begin{proof}
  Assume instead that this intersection is finite. After deleting $p$, its points still form a finite set, and all their distances from $p$ are positive.

  \begin{lean}
    classical

    intro hf

    let s := (hf.subset (show (E ∩ ball p r) \ {p} ⊆ E ∩ ball p r from sdiff_subset)).toFinset

    have hpos : ∀ q ∈ s, 0 < dist q p := by
      intro q hq

      have h := (Set.Finite.mem_toFinset _).mp hq

      exact dist_pos.mpr h.2
  \end{lean}

  Choose a common positive lower bound $\delta$ for those distances. The limit-point condition supplies a distinct point of $E$ within $\min(r,\delta)$ of $p$. It belongs to the finite set but has distance less than its lower bound, a contradiction.

  \begin{lean}
    obtain ⟨δ, hδ, hle⟩ := finite_positive_lower_bound s (fun q => dist q p) hpos
    obtain ⟨q, hq, hqp, hdist⟩ := (mem_derivedSet_iff E p).mp hp (min r δ) (lt_min hr hδ)

    have hqs : q ∈ s := (Set.Finite.mem_toFinset _).mpr
      ⟨⟨hq, hdist.trans_le (min_le_left _ _)⟩, hqp⟩

    exact (not_lt_of_ge (hle q hqs)) (hdist.trans_le (min_le_right _ _))
  \end{lean}

\end{proof}

% Lean source: Textbooks.MathematicalAnalysis1.Chapter01.MetricSpaces.open_iff_compl_derivedSet
\begin{proposition} \label{an1:thm:complement-of-open-set-is-closed}
  $E$ is open if and only if $E^{c}$ contains all its limit points, that is, $E^{c}$ is closed.

  \begin{lean}
    theorem open_iff_compl_derivedSet (E : Set X) :
        IsOpen E ↔ derivedSet Eᶜ ⊆ Eᶜ
  \end{lean}
\end{proposition}

\begin{proof}
  If $E$ is open, a point in $E$ has a ball contained in $E$. Such a ball prevents that point from being a limit point of the complement.

  \begin{lean}
    classical

    constructor

    · intro h p hp hpe

      obtain ⟨r, hr, hball⟩ := Metric.isOpen_iff.mp h p hpe
      obtain ⟨q, hq, _, hd⟩ := (mem_derivedSet_iff _ _).mp hp r hr

      exact hq (hball hd)
  \end{lean}

  Conversely, suppose the complement contains its limit points. If $p\in E$ had no ball contained in $E$, each ball would contain a point outside $E$. That point differs from $p$, making $p$ a limit point of the complement and contradicting $p\in E$.

  \begin{lean}
    · intro h
      apply Metric.isOpen_iff.mpr
      intro p hp
      by_contra hn

      have hl : p ∈ derivedSet Eᶜ := by
        rw [mem_derivedSet_iff]
        intro r hr

        have hnsub : ¬ ball p r ⊆ E := fun hs => hn ⟨r, hr, hs⟩

        obtain ⟨q, hq, hqe⟩ := Set.not_subset.mp hnsub

        refine ⟨q, hqe, ?_, hq⟩
        intro heq
        exact hqe (heq.symm ▸ hp)

      exact h hl hp
  \end{lean}

\end{proof}

% Lean source: Textbooks.MathematicalAnalysis1.Chapter01.MetricSpaces.closed_iff_derivedSet
\begin{lemma} \label{an1:lem:closed-iff-limitPoints}
  $E$ is closed if and only if $E^{\prime}\subseteq E$.

  \begin{lean}
    theorem closed_iff_derivedSet (E : Set X) :
        IsClosed E ↔ derivedSet E ⊆ E
  \end{lean}
\end{lemma}

\begin{proof}
  A set is closed precisely when its complement is open. Apply the preceding result to $E^{c}$ and take its complement twice.

  \begin{lean}
    have h := open_iff_compl_derivedSet Eᶜ

    rw [compl_compl] at h
    exact isOpen_compl_iff.symm.trans h
  \end{lean}

\end{proof}

% Lean source: Textbooks.MathematicalAnalysis1.Chapter01.MetricSpaces.perfect_iff
\begin{lemma} \label{an1:lem:perfect-iff}
  $E$ is perfect if and only if $E=E^{\prime}$.

  \begin{lean}
    theorem perfect_iff (E : Set X) : Perfect E ↔ E = derivedSet E
  \end{lean}
\end{lemma}

\begin{proof}
  Closedness is equivalent to $E^{\prime}\subseteq E$, by the preceding result. Every point being a limit point says $E\subseteq E^{\prime}$. Together these are equality.

  \begin{lean}
    rw [perfect_def]
    change (IsClosed E ∧ E ⊆ derivedSet E) ↔ E = derivedSet E
    rw [closed_iff_derivedSet]
    exact ⟨fun h => subset_antisymm h.2 h.1, fun h => ⟨h.symm.subset, h.subset⟩⟩
  \end{lean}
\end{proof}

% Lean source: Textbooks.MathematicalAnalysis1.Chapter01.MetricSpaces.open_union
\begin{proposition} \label{an1:thm:arbitrary-union-of-open-sets-is-open}
  For any indexed family of open sets $(E_i)$, the union $\bigcup_i E_i$ is open.

  \begin{lean}
    theorem open_union {ι : Type*} (E : ι → Set X)
        (h : ∀ i, IsOpen (E i)) : IsOpen (⋃ i, E i)
  \end{lean}
\end{proposition}

\begin{proof}
  A point in the union belongs to one member $E_i$. Choose a ball contained in that member; the same ball is contained in the union. This argument also covers an empty family, since then there is no point to consider.

  \begin{lean}
    apply Metric.isOpen_iff.mpr
    intro p hp

    obtain ⟨i, hi⟩ := mem_iUnion.mp hp
    obtain ⟨r, hr, hs⟩ := Metric.isOpen_iff.mp (h i) p hi

    exact ⟨r, hr, fun q hq => mem_iUnion.mpr ⟨i, hs hq⟩⟩
  \end{lean}

\end{proof}

% Lean source: Textbooks.MathematicalAnalysis1.Chapter01.MetricSpaces.closed_intersection
\begin{corollary} \label{an1:cor:arbitrary-intersection-of-closed-sets-is-closed}
  For any indexed family of closed sets $(E_i)$, the intersection $\bigcap_i E_i$ is closed.

  \begin{lean}
    theorem closed_intersection {ι : Type*} (E : ι → Set X)
        (h : ∀ i, IsClosed (E i)) : IsClosed (⋂ i, E i)
  \end{lean}
\end{corollary}

\begin{proof}
  A limit point of the intersection is a limit point of each member: every nearby witness in the intersection belongs to that member. Each member is closed, so it contains the point; consequently the intersection contains it.

  \begin{lean}
    apply (closed_iff_derivedSet _).mpr
    intro p hp
    apply mem_iInter.mpr
    intro i
    apply (closed_iff_derivedSet _).mp (h i)
    rw [mem_derivedSet_iff]
    intro r hr

    obtain ⟨q, hq, hn, hd⟩ := (mem_derivedSet_iff _ _).mp hp r hr

    exact ⟨q, mem_iInter.mp hq i, hn, hd⟩
  \end{lean}

\end{proof}

% Lean source: Textbooks.MathematicalAnalysis1.Chapter01.MetricSpaces.open_intersection
\begin{lemma} \label{an1:lem:open-intersection}
  The intersection of two open sets is open.

  \begin{lean}
    theorem open_intersection (E F : Set X) (hE : IsOpen E) (hF : IsOpen F) :
        IsOpen (E ∩ F)
  \end{lean}
\end{lemma}

\begin{proof}
  At a point of the intersection, choose one positive radius for each set. The smaller radius gives a ball lying in both sets.

  \begin{lean}
    apply Metric.isOpen_iff.mpr
    intro p hp

    obtain ⟨r, hr, hEr⟩ := Metric.isOpen_iff.mp hE p hp.1
    obtain ⟨s, hs, hFs⟩ := Metric.isOpen_iff.mp hF p hp.2

    refine ⟨min r s, lt_min hr hs, ?_⟩
    intro q hq
    change dist q p < min r s at hq
    exact ⟨hEr (hq.trans_le (min_le_left _ _)), hFs (hq.trans_le (min_le_right _ _))⟩
  \end{lean}

\end{proof}

% Lean source: Textbooks.MathematicalAnalysis1.Chapter01.MetricSpaces.open_finite_intersection
\begin{proposition} \label{an1:thm:finite-intersection-of-open-sets-is-open}
  A finite intersection of open sets is open. In particular this applies to $\bigcap_{j=0}^{n-1}E_j$ for every $n\in\N$.

  \begin{lean}
    theorem open_finite_intersection {ι : Type*} (s : Finset ι)
        (E : ι → Set X) (hE : ∀ i ∈ s, IsOpen (E i)) :
        IsOpen (⋂ i ∈ s, E i)
  \end{lean}
\end{proposition}

\begin{proof}
  For no indices the intersection is $X$, which contains a unit ball about each of its points. Adjoining one index replaces an already open finite intersection by its intersection with one more open set, so induction completes the proof.

  \begin{lean}
    classical

    induction s using Finset.induction_on with
    | empty =>
      simp only [Finset.notMem_empty, iInter_of_empty, iInter_univ]
      exact Metric.isOpen_iff.mpr (fun p _ => ⟨1, zero_lt_one, subset_univ _⟩)

    | @insert i s hi ih =>
      rw [Finset.set_biInter_insert]
      exact open_intersection _ _ (hE i (Finset.mem_insert_self i s))
        (ih (fun j hj => hE j (Finset.mem_insert_of_mem hj)))
  \end{lean}

\end{proof}

% Lean source: Textbooks.MathematicalAnalysis1.Chapter01.MetricSpaces.closed_finite_union
\begin{corollary} \label{an1:cor:finite-union-of-closed-sets-is-closed}
  A finite union of closed sets is closed, including the union over no indices.

  \begin{lean}
    theorem closed_finite_union {ι : Type*} (s : Finset ι)
        (E : ι → Set X) (hE : ∀ i ∈ s, IsClosed (E i)) :
        IsClosed (⋃ i ∈ s, E i)
  \end{lean}
\end{corollary}

\begin{proof}
  The complements are open, so their finite intersection is open. By De Morgan\textquotesingle s identity, this is the complement of the desired union; hence that union is closed.

  \begin{lean}
    have ho := open_finite_intersection s (fun i => (E i)ᶜ)
      (fun i hi => (hE i hi).isOpen_compl)

    have heq : (⋃ i ∈ s, E i)ᶜ = ⋂ i ∈ s, (E i)ᶜ := by
      ext p
      simp only [mem_compl_iff, mem_iUnion, mem_iInter, not_exists]

    exact isOpen_compl_iff.mp (heq ▸ ho)
  \end{lean}

\end{proof}

% Lean source: Textbooks.MathematicalAnalysis1.Chapter01.MetricSpaces.closure_eq_union_derivedSet
\begin{lemma} \label{an1:lem:closure-eq-union-limitPoints}
  $\overline E=E\cup E^{\prime}$, relating closure to the derived set.

  \begin{lean}
    theorem closure_eq_union_derivedSet (E : Set X) :
        closure E = E ∪ derivedSet E
  \end{lean}
\end{lemma}

\begin{proof}
  If $p\in\overline E$, every ball about $p$ meets $E$. Either $p\in E$, or each such witness differs from $p$, making $p$ a limit point.

  \begin{lean}
    classical

    ext p
    constructor

    · intro hp
      by_cases he : p ∈ E

      · exact Or.inl he
      · right
        rw [mem_derivedSet_iff]
        intro r hr

        obtain ⟨q, hq, hd⟩ := Metric.mem_closure_iff.mp hp r hr

        exact ⟨q, hq, fun h => he (h ▸ hq), by rwa [dist_comm]⟩
  \end{lean}

  Conversely, a point of $E$ itself witnesses that every ball meets $E$. For a limit point, use the distinct witness in each deleted ball. Both cases give membership in the closure.

  \begin{lean}
    · intro hp
      apply Metric.mem_closure_iff.mpr
      intro r hr
      rcases hp with hp | hp

      · exact ⟨p, hp, by simpa only [dist_self] using hr⟩

      · obtain ⟨q, hq, _, hd⟩ := (mem_derivedSet_iff _ _).mp hp r hr

        exact ⟨q, hq, by rwa [dist_comm]⟩
  \end{lean}

\end{proof}

% Lean source: Textbooks.MathematicalAnalysis1.Chapter01.MetricSpaces.subset_closure
\begin{lemma} \label{an1:lem:subset-closure}
  Every set is contained in its closure.

  \begin{lean}
    theorem subset_closure (E : Set X) : E ⊆ closure E
  \end{lean}
\end{lemma}

\begin{proof}
  By the preceding characterization, every point of $E$ belongs to the left member of the union $E\cup E^{\prime}$.

  \begin{lean}
    rw [closure_eq_union_derivedSet]
    exact fun _ hp => Or.inl hp
  \end{lean}

\end{proof}

% Lean source: Textbooks.MathematicalAnalysis1.Chapter01.MetricSpaces.closure_closed
\begin{lemma} \label{an1:lem:closure-closed}
  The closure of any set is closed.

  \begin{lean}
    theorem closure_closed (E : Set X) : IsClosed (closure E)
  \end{lean}
\end{lemma}

\begin{proof}
  Suppose $p$ is a limit point of $\overline E$ and $r>0$. Choose $q\in\overline E$ within $r/2$ of $p$, then $s\in E$ within $r/2$ of $q$. The triangle inequality puts $s$ within $r$ of $p$, proving $p\in\overline E$.

  \begin{lean}
    apply (closed_iff_derivedSet _).mpr
    intro p hp
    apply Metric.mem_closure_iff.mpr
    intro r hr

    obtain ⟨q, hq, _, hdq⟩ := (mem_derivedSet_iff _ _).mp hp (r / 2) (half_pos hr)
    obtain ⟨s, hs, hds⟩ := Metric.mem_closure_iff.mp hq (r / 2) (half_pos hr)

    refine ⟨s, hs, ?_⟩

    have ht := dist_triangle p q s

    rw [dist_comm q p] at hdq
    linarith
  \end{lean}

\end{proof}

% Lean source: Textbooks.MathematicalAnalysis1.Chapter01.MetricSpaces.closure_properties
\begin{proposition} \label{an1:thm:properties-of-closure}
  For every $E\subseteq X$: $\overline E$ is closed; $E=\overline E$ if and only if $E$ is closed; and every closed set $F$ containing $E$ also contains $\overline E$.

  \begin{lean}
    theorem closure_properties (E : Set X) :
        IsClosed (closure E) ∧ (E = closure E ↔ IsClosed E) ∧
        (∀ F : Set X, IsClosed F → E ⊆ F → closure E ⊆ F)
  \end{lean}
\end{proposition}

\begin{proof}
  We have just proved that $\overline E$ is closed. Therefore $E=\overline E$ implies that $E$ is closed. Conversely, closedness gives $E^{\prime}\subseteq E$, so $E\cup E^{\prime}=E$.

  \begin{lean}
    refine ⟨closure_closed E, ?_, ?_⟩

    · constructor
      · intro h
        rw [h]
        exact closure_closed E

      · intro h
        rw [closure_eq_union_derivedSet]
        exact (union_eq_left.mpr ((closed_iff_derivedSet E).mp h)).symm
  \end{lean}

  Let $F$ be closed with $E\subseteq F$. A point of $\overline E$ either lies in $E$, hence in $F$, or is a limit point of $E$. In the latter case the same nearby witnesses make it a limit point of $F$, which belongs to $F$ by closedness.

  \begin{lean}
    · intro F hF hEF p hp
      rw [closure_eq_union_derivedSet] at hp
      rcases hp with hp | hp

      · exact hEF hp
      · apply (closed_iff_derivedSet F).mp hF
        rw [mem_derivedSet_iff]
        intro r hr

        obtain ⟨q, hq, hn, hd⟩ := (mem_derivedSet_iff _ _).mp hp r hr

        exact ⟨q, hEF hq, hn, hd⟩
  \end{lean}

\end{proof}

% Lean source: Textbooks.MathematicalAnalysis1.Chapter01.MetricSpaces.supremum_in_closure
\begin{proposition} \label{an1:thm:supremum-in-closure}
  If $E\subseteq\R$ is nonempty and bounded above, then $\sup E\in\overline E$. If $E$ is also closed, then $\sup E\in E$.

  \begin{lean}
    theorem supremum_in_closure (E : Set ℝ) (hne : E.Nonempty) (hb : BddAbove E) :
        sSup E ∈ closure E ∧ (IsClosed E → sSup E ∈ E)
  \end{lean}
\end{proposition}

\begin{proof}
  Write $c=\sup E$. For any $r>0$, the smaller number $c-r$ is not an upper bound, so choose $q\in E$ with $c-r<q\leq c$. Then $|c-q|<r$, proving that every ball about $c$ meets $E$.

  \begin{lean}
    have hc : sSup E ∈ closure E := by
      apply Metric.mem_closure_iff.mpr
      intro r hr

      obtain ⟨q, hq, hlt⟩ := exists_lt_of_lt_csSup hne (sub_lt_self (sSup E) hr)

      refine ⟨q, hq, ?_⟩
      rw [Real.dist_eq, abs_of_nonneg (sub_nonneg.mpr (le_csSup hb hq))]
      linarith
  \end{lean}

  When $E$ is closed, its closure equals $E$, so the established closure membership is membership in $E$ itself.

  \begin{lean}
    exact ⟨hc, fun h => by
      have heq := (closure_properties E).2.1.mpr h

      rw [← heq] at hc
      exact hc⟩
  \end{lean}

\end{proof}

% Lean source: Textbooks.MathematicalAnalysis1.Chapter01.MetricSpaces.isolated_eq_diff
\begin{proposition} \label{an1:thm:isolated-points-are-self-minus-limit-points}
  $\operatorname{Iso}(E)=E\setminus E^{\prime}$.

  \begin{lean}
    theorem isolated_eq_diff (E : Set X) : isolatedPoints E = E \ derivedSet E
  \end{lean}
\end{proposition}

\begin{proof}
  If $E\cap B(p,r)=\{p\}$, then $p\in E$ and no distinct point of $E$ is in that ball. This contradicts the limit-point condition, so $p\notin E^{\prime}$.

  \begin{lean}
    classical

    ext p
    constructor

    · rintro ⟨r, hr, heq⟩

      have hp : p ∈ E ∩ ball p r := heq.symm ▸ (mem_singleton p)

      refine ⟨hp.1, ?_⟩
      intro hl

      obtain ⟨q, hq, hn, hd⟩ := (mem_derivedSet_iff _ _).mp hl r hr

      have hmem : q ∈ E ∩ ball p r := ⟨hq, hd⟩

      rw [heq, mem_singleton_iff] at hmem
      exact hn hmem
  \end{lean}

  Conversely, if $p\in E$ but $p\notin E^{\prime}$, some positive-radius ball contains no distinct point of $E$. Its intersection with $E$ is therefore exactly $\{p\}$.

  \begin{lean}
    · rintro ⟨hp, hn⟩

      have hex : ∃ r > 0, ∀ q ∈ E, dist q p < r → q = p := by
        by_contra h
        push Not at h
        apply hn
        rw [mem_derivedSet_iff]
        intro r hr

        obtain ⟨q, hq, hd, hne⟩ := h r hr

        exact ⟨q, hq, hne, hd⟩

      obtain ⟨r, hr, h⟩ := hex

      refine ⟨r, hr, ?_⟩
      ext q
      constructor

      · intro hq
        exact h q hq.1 hq.2

      · intro hq
        rw [mem_singleton_iff] at hq
        subst q
        exact ⟨hp, by simpa only [mem_ball, dist_self] using hr⟩
  \end{lean}

\end{proof}

% Lean source: Textbooks.MathematicalAnalysis1.Chapter01.MetricSpaces.compl_interior_eq_closure_compl
\begin{lemma} \label{an1:lem:complement-of-interior-is-closure-of-complement}
  $(E^{\circ})^{c}=\overline{E^{c}}$.

  \begin{lean}
    theorem compl_interior_eq_closure_compl (E : Set X) :
        (interior E)ᶜ = closure Eᶜ
  \end{lean}
\end{lemma}

\begin{proof}
  Failure to be interior means that no positive-radius ball lies in $E$. Thus every such ball contains a point of $E^{c}$, giving membership in its closure. Conversely, if every ball meets $E^{c}$, no ball can be contained in $E$.

  \begin{lean}
    classical

    ext p
    rw [mem_compl_iff, mem_interior_iff_mem_nhds, Metric.mem_nhds_iff,
      Metric.mem_closure_iff]
    constructor

    · intro h r hr

      have hn : ¬ ball p r ⊆ E := fun hs => h ⟨r, hr, hs⟩

      obtain ⟨q, hq, hqe⟩ := Set.not_subset.mp hn

      exact ⟨q, hqe, by rwa [dist_comm]⟩

    · intro h ⟨r, hr, hs⟩

      obtain ⟨q, hq, hd⟩ := h r hr

      exact hq (hs (by simpa only [mem_ball, dist_comm] using hd))
  \end{lean}

\end{proof}

\begin{definition} \label{an1:def:boundary}
  The \textbf{boundary} is $\partial E=\overline E\setminus E^{\circ}$.
  The next lemma identifies it with $\overline E\cap\overline{E^{c}}$.
\end{definition}

% Lean source: Textbooks.MathematicalAnalysis1.Chapter01.MetricSpaces.boundary_eq
\begin{lemma} \label{an1:lem:boundary-eq}
  $\partial E=\overline E\cap\overline{E^{c}}$.

  \begin{lean}
    theorem boundary_eq (E : Set X) :
        frontier E = closure E ∩ closure Eᶜ
  \end{lean}
\end{lemma}

\begin{proof}
  Write the difference as an intersection with the complement of the interior. The preceding lemma identifies this complement with the closure of the complement.

  \begin{lean}
    rw [frontier, sdiff_eq_compl_inter, compl_interior_eq_closure_compl, inter_comm]
  \end{lean}

\end{proof}

% Lean source: Textbooks.MathematicalAnalysis1.Chapter01.EuclideanSpaces.zero_dimensional
\begin{lemma} \label{an1:lem:zero-dimensional}
  Any two points of $\R^{0}$ are equal, and every punctured ball in $\R^{0}$ is empty.

  \begin{lean}
    theorem zero_dimensional (p q : EuclideanSpace ℝ (Fin 0)) (r : ℝ) :
        p = q ∧ ball p r \ {p} = ∅
  \end{lean}
\end{lemma}

\begin{proof}
  There are no coordinates on which two empty tuples can differ, so they are equal. Removing the center consequently removes every possible point of a ball.

  \begin{lean}
    have heq : ∀ x y : EuclideanSpace ℝ (Fin 0), x = y := by
      intro x y
      ext j
      exact Fin.elim0 j

    refine ⟨heq p q, ?_⟩
    apply Set.eq_empty_iff_forall_notMem.mpr
    intro x hx
    exact hx.2 (heq x p)
  \end{lean}

\end{proof}

% Lean source: Textbooks.MathematicalAnalysis1.Chapter01.EuclideanSpaces.euclidean_punctured_ball
\begin{lemma} \label{an1:lem:deleted-neighborhood-in-Rn-is-non-empty}
  For $n\in\N^{+}$, $p\in\R^{n}$, and $r>0$, the punctured ball $B(p,r)\setminus\{p\}$ is nonempty.

  \begin{lean}
    theorem euclidean_punctured_ball (n : ℕ+)
        (p : EuclideanSpace ℝ (Fin n)) (r : ℝ) (hr : 0 < r) :
        ∃ q : EuclideanSpace ℝ (Fin n), q ≠ p ∧ dist q p < r
  \end{lean}
\end{lemma}

\begin{proof}
  Since $n>0$, the coordinate vector $e_0$ exists. Put $v=(r/2)e_0$, so $\|v\|=r/2>0$. The point $p+v$ differs from $p$ and is at distance $r/2<r$ from it.

  \begin{lean}
    let v : EuclideanSpace ℝ (Fin n) := EuclideanSpace.single ⟨0, n.pos⟩ (r / 2)

    have hv : ‖v‖ = r / 2 := by
      rw [PiLp.norm_single, Real.norm_eq_abs, abs_of_pos (half_pos hr)]

    refine ⟨p + v, ?_, ?_⟩

    · intro heq

      have hz : v = 0 := add_left_cancel (heq.trans (add_zero p).symm)

      rw [hz, norm_zero] at hv
      linarith

    · rw [dist_eq_norm, add_sub_cancel_left, hv]
      linarith
  \end{lean}

\end{proof}

% Lean source: Textbooks.MathematicalAnalysis1.Chapter01.MetricSpaces.isolated_limit_compl
\begin{lemma} \label{an1:lem:isolated-limit-compl}
  If every positive-radius deleted ball in $X$ is nonempty, then $\operatorname{Iso}(E)\subseteq(E^{c})^{\prime}$.

  \begin{lean}
    theorem isolated_limit_compl (E : Set X)
        (hX : ∀ p : X, ∀ r > 0, ∃ q : X, q ≠ p ∧ dist q p < r) :
        isolatedPoints E ⊆ derivedSet Eᶜ
  \end{lean}
\end{lemma}

\begin{proof}
  For an isolated point $p$, choose $r>0$ with $E\cap B(p,r)=\{p\}$. For any $s>0$, take a distinct point within $\min(r,s)$ of $p$. It must lie outside $E$, and therefore witnesses that $p$ is a limit point of $E^{c}$.

  \begin{lean}
    intro p hp

    obtain ⟨r, hr, heq⟩ := hp

    rw [mem_derivedSet_iff]
    intro s hs

    obtain ⟨q, hne, hd⟩ := hX p (min r s) (lt_min hr hs)

    refine ⟨q, ?_, hne, hd.trans_le (min_le_right _ _)⟩
    intro hq

    have hm : q ∈ E ∩ ball p r := ⟨hq, hd.trans_le (min_le_left _ _)⟩

    rw [heq, mem_singleton_iff] at hm
    exact hne hm
  \end{lean}

\end{proof}

% Lean source: Textbooks.MathematicalAnalysis1.Chapter01.EuclideanSpaces.euclidean_boundary
\begin{proposition} \label{an1:thm:properties-of-boundary-in-rn}
  For $n\in\N^{+}$ and $E\subseteq\R^{n}$, $\partial E=\operatorname{Iso}(E)\cup\operatorname{Iso}(E^{c})\cup(E^{\prime}\cap(E^{c})^{\prime})$.

  \begin{lean}
    theorem euclidean_boundary (n : ℕ+)
        (E : Set (EuclideanSpace ℝ (Fin n))) :
        frontier E = isolatedPoints E ∪ isolatedPoints Eᶜ ∪
          (derivedSet E ∩ derivedSet Eᶜ)
  \end{lean}
\end{proposition}

\begin{proof}
  An isolated point of either side is a limit point of the other, by the preceding two results. Expand the two closures in $\partial E$ as the set together with its limit points. Separate points which fail one of the two limit-point conditions from those satisfying both. The former are exactly the isolated points of one side; the latter form $E^{\prime}\cap(E^{c})^{\prime}$.

  \begin{lean}
    have hE := isolated_limit_compl E (euclidean_punctured_ball n)

    have hEc := isolated_limit_compl Eᶜ (euclidean_punctured_ball n)

    rw [compl_compl] at hEc
    rw [isolated_eq_diff] at hE hEc
    rw [frontier, sdiff_eq_compl_inter,
      compl_interior_eq_closure_compl, closure_eq_union_derivedSet,
      closure_eq_union_derivedSet, isolated_eq_diff, isolated_eq_diff]
    ext p

    have h1 := @hE p

    have h2 := @hEc p

    simp only [mem_inter_iff, mem_union, mem_sdiff, mem_compl_iff] at *
    tauto
  \end{lean}

\end{proof}

% Lean source: Textbooks.MathematicalAnalysis1.Chapter01.MetricSpaces.subspace_metric
\begin{proposition} \label{an1:lem:subspace-metric}
  For $Y\subseteq X$, restricting $d$ to $Y\times Y$ gives a metric on $Y$, with the same distances.

  \begin{lean}
    theorem subspace_metric (Y : Set X) (p q z : Y) :
        dist p q = dist p.val q.val ∧ 0 ≤ dist p q ∧
        (dist p q = 0 ↔ p = q) ∧ dist p q = dist q p ∧
        dist p z ≤ dist p q + dist q z
  \end{lean}
\end{proposition}

\begin{proof}
  Nonnegativity, symmetry, and the triangle inequality are inherited from $X$. Distance zero means equality of the underlying points, which is also equality as points of $Y$.

  \begin{lean}
    exact ⟨rfl, dist_nonneg, dist_eq_zero, dist_comm _ _, dist_triangle _ _ _⟩
  \end{lean}

\end{proof}

The resulting space is called the \textbf{subspace} $(Y,d)$. We continue to write $d(p,q)$ for its distance.

\begin{definition} \label{an1:def:relative-open-closed-sets}
  For $E\subseteq Y\subseteq X$, $E$ is \textbf{open relative to} $Y$ when it is open as a subset of the subspace $Y$. Equivalently, each $p\in E$ has $r>0$ with $Y\cap B(p,r)\subseteq E$.
  It is \textbf{closed relative to} $Y$ when it is closed in that subspace. The next lemma characterizes this by $E^{\prime}\cap Y\subseteq E$.
\end{definition}

% Lean source: Textbooks.MathematicalAnalysis1.Chapter01.MetricSpaces.relative_closed
\begin{lemma} \label{an1:lem:relative-closed}
  For $E\subseteq Y$, $E$ is closed in $Y$ if and only if $E^{\prime}\cap Y\subseteq E$.

  \begin{lean}
    theorem relative_closed (E Y : Set X) (hEY : E ⊆ Y) :
        IsClosed ((Subtype.val : Y → X) ⁻¹' E) ↔ derivedSet E ∩ Y ⊆ E
  \end{lean}
\end{lemma}

\begin{proof}
  A point of $Y$ that is an ambient limit point of $E$ is a subspace limit point: every witness lies in $E\subseteq Y$, and its distance is unchanged.

  \begin{lean}
    rw [closed_iff_derivedSet]
    constructor
  \end{lean}

  Conversely, each subspace witness has an underlying point in $E$ with the same positive distance. Thus a subspace limit point is an ambient limit point lying in $Y$, and the stated condition puts it in $E$.

  \begin{lean}
    · intro h p hp
      apply h (show (⟨p, hp.2⟩ : Y) ∈ derivedSet (Subtype.val ⁻¹' E) from ?_)
      rw [mem_derivedSet_iff]
      intro r hr

      obtain ⟨q, hq, hne, hd⟩ := (mem_derivedSet_iff _ _).mp hp.1 r hr

      refine ⟨⟨q, hEY hq⟩, hq, ?_, hd⟩
      intro heq
      exact hne (congrArg Subtype.val heq)

    · intro h p hp
      apply h
      refine ⟨?_, p.property⟩
      rw [mem_derivedSet_iff]
      intro r hr

      obtain ⟨q, hq, hne, hd⟩ := (mem_derivedSet_iff _ _).mp hp r hr

      exact ⟨q.val, hq, fun heq => hne (Subtype.ext heq), hd⟩
  \end{lean}

\end{proof}

% Lean source: Textbooks.MathematicalAnalysis1.Chapter01.MetricSpaces.relative_open
\begin{proposition} \label{an1:thm:open-set-relative-to-subspace}
  For $E\subseteq Y\subseteq X$, $E$ is open in $Y$ if and only if $E=V\cap Y$ for some open $V\subseteq X$.

  \begin{lean}
    theorem relative_open (E Y : Set X) (hEY : E ⊆ Y) :
        IsOpen ((Subtype.val : Y → X) ⁻¹' E) ↔
          ∃ V : Set X, IsOpen V ∧ E = V ∩ Y
  \end{lean}
\end{proposition}

\begin{proof}
  Suppose $E$ is open in $Y$. For each $p\in E$, choose $r_p>0$ with $Y\cap B(p,r_p)\subseteq E$.

  \begin{lean}
    classical

    constructor

    · intro h

      have hex : ∀ p : E, ∃ r > 0, Y ∩ ball p.val r ⊆ E := by
        intro p

        obtain ⟨r, hr, hs⟩ := Metric.isOpen_iff.mp h ⟨p.val, hEY p.property⟩ p.property

        exact ⟨r, hr, fun q hq =>
          hs (show (⟨q, hq.1⟩ : Y) ∈ ball ⟨p.val, hEY p.property⟩ r from hq.2)⟩
  \end{lean}

  Let $V$ be the union of these ambient balls. It is open. Each point of $E$ lies in its own ball and in $Y$; conversely, a point in $V\cap Y$ lies in one chosen ball, whose intersection with $Y$ lies in $E$. Hence $E=V\cap Y$.

  \begin{lean}
    choose r hr hs using hex

    refine ⟨⋃ p : E, ball p.val (r p), open_union _ (fun p => ball_open _ _), ?_⟩
    ext q
    constructor

    · intro hq
      exact ⟨mem_iUnion.mpr ⟨⟨q, hq⟩, by simpa only [mem_ball, dist_self] using hr ⟨q, hq⟩⟩, hEY hq⟩

    · rintro ⟨hq, hy⟩

      obtain ⟨p, hp⟩ := mem_iUnion.mp hq

      exact hs p ⟨hy, hp⟩
  \end{lean}

  Conversely, choose an ambient ball inside $V$ about each point of $E=V\cap Y$. Restricting this ball to $Y$ produces a subspace ball inside $E$.

  \begin{lean}
    · rintro ⟨V, hV, heq⟩
      apply Metric.isOpen_iff.mpr
      intro p hp

      have hpV : p.val ∈ V := (show p.val ∈ V ∩ Y from heq ▸ hp).1

      obtain ⟨r, hr, hs⟩ := Metric.isOpen_iff.mp hV p.val hpV

      refine ⟨r, hr, ?_⟩
      intro q hq
      change q.val ∈ E
      rw [heq]
      exact ⟨hs hq, q.property⟩
  \end{lean}

\end{proof}
